% Generated draft: requires mathematical and source review.
\documentclass[11pt]{article}
\usepackage[margin=25mm]{geometry}
\usepackage{fontspec}
\setmainfont{Latin Modern Roman}
\usepackage{amsmath,amssymb,mathtools,unicode-math}
\setmathfont{Latin Modern Math}
\usepackage{enumitem,needspace,xurl,hyperref}
\setlength{\parindent}{0pt}
\setlength{\parskip}{5pt}
\setlength{\emergencystretch}{3em}
\newcommand{\N}{\mathbb{N}}
\newcommand{\Z}{\mathbb{Z}}
\newcommand{\Q}{\mathbb{Q}}
\newcommand{\R}{\mathbb{R}}
\newcommand{\C}{\mathbb{C}}
\newcommand{\F}{\mathbb{F}}
\newcommand{\A}{\mathbb{A}}
\newcommand{\PP}{\mathbb{P}}
\newcommand{\E}{\mathbb{E}}
\newcommand{\eps}{\varepsilon}
\newcommand{\dd}{\,\mathrm{d}}
\begin{document}
\section*{UnsolvedMath ID: 30002875}
\textbf{Source problem code:} catalogue entry
\subsection{Definitions and mathematical statement}
Interpret \(G\) as a connected simple simply connected algebraic group over an algebraically closed field \(k\) of characteristic \(p>0\). Choose a maximal torus \(T\) and positive roots, and let \(B\) be the Borel subgroup corresponding to the negative roots. Write \(X(T)\) for the weight lattice, \(X^+\) for its dominant weights, and \(\rho\) for half the sum of the positive roots. For \(\lambda\in X^+\), let \(L(\lambda)\) denote the simple rational \(G\)-module of highest weight \(\lambda\), and let \(\nabla(\lambda)=\operatorname{ind}_B^G k_\lambda\) denote the induced module of highest weight \(\lambda\), where \(k_\lambda\) is the one-dimensional \(B\)-module defined by \(\lambda\).

For an integer \(r\geq 1\), define the set of \(p^r\)-restricted dominant weights by
\[
X_r=\{\lambda\in X^+:0\leq\langle\lambda,\alpha^\vee\rangle<p^r\text{ for every simple coroot }\alpha^\vee\}.
\]
Every \(\lambda\in X^+\) has a unique expression \(\lambda=\lambda_0+p^r\lambda_1\), with \(\lambda_0\in X_r\) and \(\lambda_1\in X^+\). Set
\[
\nabla^{(p,r)}(\lambda)=L(\lambda_0)\otimes\nabla(\lambda_1)^{(r)},
\qquad
\mathrm{St}_r=L((p^r-1)\rho).
\]
Here \(V^{(r)}\) denotes the \(r\)-th Frobenius twist, using a standard split Frobenius morphism of \(G\).

A finite-dimensional rational \(G\)-module \(M\) has a good filtration if it admits a finite filtration by \(G\)-submodules
\[
0=M_0\subset M_1\subset\cdots\subset M_n=M
\]
whose successive quotients are modules \(\nabla(\mu_i)\) with \(\mu_i\in X^+\). It has a good \((p,r)\)-filtration if it admits such a filtration with successive quotients \(\nabla^{(p,r)}(\mu_i)\). The empty filtration is allowed when \(M=0\).

The supplied formulation of Donkin's Good Filtration Conjecture asks whether, for every such \(k\) and \(G\), every integer \(r\geq1\), and every finite-dimensional rational \(G\)-module \(M\),
\[
M\text{ has a good }(p,r)\text{-filtration}
\quad\Longleftrightarrow\quad
\mathrm{St}_r\otimes_k M\text{ has a good filtration}.
\]
The algebraic-group interpretation, the positive-integer range of \(r\), and the filtration definitions above make explicit standard conventions omitted from the supplied statement; their agreement with the intended source requires review. The supplied statement imposes no further restriction on \(p\). Its accompanying literature assessment concerns composition-factor multiplicities of induced modules, which is a different stated target and does not by itself establish the status of this equivalence.
\subsection{Short English statement}
The supplied conjecture asserts that a finite-dimensional rational module has a good \((p,r)\)-filtration exactly when tensoring it with the \(r\)-th Steinberg module produces a module with a good filtration. The source labels it open, but its accompanying literature assessment addresses a different problem; the formulation and status require human review.
\subsection{Sources}
\begin{itemize}
\item Supplied local UnsolvedMath registry record, permanent ID \(30002875\), source problem code \texttt{catalogue entry}, title \emph{Donkin's Good Filtration Conjecture}; registry context: \texttt{data/unsolvedmath/problems.json}. The record attributes the conjecture to Donkin, lists no individual proposer, and records proposed year \(2015\).
\item Supplied report citation: \emph{Cohomology of Finite Groups: Interactions and Applications} (2016), Oberwolfach Reports, pp.\ 1317--1384, DOI \texttt{10.4171/owr/2015/24}. \url{https://doi.org/10.4171/owr/2015/24}. The background describes Oberwolfach Reports as records of mathematical workshops and their open-problem discussions. The supplied citation year \(2016\) and the report identifier \(2015/24\) are retained without reconciliation; no problem-specific page or original Donkin reference is supplied.
\item Supplied literature source: Williamson, \emph{Algebraic representations and constructible sheaves}. \url{https://arxiv.org/abs/1610.06261}.
\item Supplied literature source: \emph{Oberwolfach Report 2015/24}. \url{https://doi.org/10.4171/owr/2015/24}.
\item Literature checked at: \texttt{2026-08-21}. This is supplied dated curation derived from the supplied OWR curation, not an independently verified assessment. It labels the problem open and says that the general composition-factor multiplicity problem for induced modules in modular reductive-group representation theory remains unresolved. That assessment concerns a different stated target from the filtration equivalence. Time-sensitive status and cited claims require independent verification.
\item The supplied statement status is \texttt{corrected\_verified}; its verification note says that mathematical content was preserved while notation, typography, or omitted definitions were repaired for publication. These are registry assertions, not independent mathematical verification of this draft.
\end{itemize}
\end{document}
